\chapter{PENDAHULUAN}

\section{Latar Belakang}

Restoran Nasi Padang merupakan salah satu jenis rumah makan yang mudah ditemukan di Indonesia. Salah satu karakteristiknya adalah tersedianya berbagai jenis lauk yang dapat dipilih pelanggan, seperti rendang, perkedel, telur dadar, sayur nangka, daun singkong, telur kari, dan sambal. Pada pelayanan konvensional, pelanggan menyampaikan lauk yang diinginkan kepada staf, kemudian staf mengambilkan lauk tersebut untuk pelanggan. Setelah selesai makan, pelanggan menghampiri staf pada bagian pembayaran dan menyebutkan satu per satu jenis serta jumlah lauk yang telah dikonsumsi berdasarkan ingatan pelanggan. Staf kemudian menghitung total harga dari lauk yang disebutkan dan menyampaikan jumlah yang harus dibayarkan kepada pelanggan.

Dalam kondisi jumlah pelanggan yang tinggi, proses pengambilan lauk tersebut dapat menyebabkan pelanggan harus menunggu ketersediaan staf sebelum dapat dilayani. Pada saat yang sama, staf perlu membagi perhatian antara melayani pelanggan \textit{dine-in} dan menangani pesanan lain, termasuk pesanan \textit{take-away}. Kondisi ini dapat menimbulkan antrean pada proses pelayanan dan mengurangi efisiensi kerja staf.

Untuk mengatasi permasalahan tersebut, penelitian ini mengusulkan alternatif pelayanan berupa \textit{self-service} untuk pelanggan \textit{dine-in}. Pada mekanisme ini, pelanggan dapat mengambil dan memilih sendiri lauk yang diinginkan dari lauk yang tersedia, kemudian membawa piring yang telah dipilih ke proses \textit{checkout}. Dengan demikian, pelanggan tidak perlu menunggu staf untuk mengambilkan lauk satu per satu. Penerapan mekanisme \textit{self-service} juga memungkinkan staf untuk mengurangi aktivitas pelayanan pengambilan lauk \textit{dine-in} sehingga dapat lebih berfokus pada pelayanan lain, khususnya pemenuhan pesanan \textit{take-away}.

Agar mekanisme tersebut dapat diterapkan tanpa mengharuskan staf memeriksa
kembali lauk yang dipilih pelanggan secara \textit{self-service}, diperlukan
suatu sistem yang mampu mengenali isi piring secara otomatis. Kebutuhan
tersebut merupakan salah satu permasalahan yang dapat ditangani menggunakan
\textit{Computer Vision}. Pemanfaatan \textit{Computer Vision} dalam analisis
makanan telah mencakup berbagai tugas seperti pengenalan, deteksi, dan
segmentasi makanan \cite{min2019}.

Dalam beberapa kasus lauk yang terdapat pada satu piring, pengenalan kelas
saja belum cukup karena sistem juga perlu membedakan setiap lauk sebagai
objek yang terpisah untuk menghitung jumlahnya. Oleh karena itu, penelitian
ini menggunakan \textit{instance segmentation}. \textit{Instance segmentation}
menghasilkan informasi kelas sekaligus \textit{mask} untuk setiap instans
objek sehingga beberapa objek dari kelas yang sama tetap dapat
direpresentasikan secara terpisah \cite{he2017maskrcnn}.

Kemampuan tersebut sesuai dengan karakteristik objek makanan yang dapat
memiliki bentuk tidak beraturan dan saling menutupi serta dengan kebutuhan
sistem \textit{checkout} untuk menghitung setiap instans secara individual
\cite{guo2023,nguyen2022sibnet,nguyen2023foodmask}. Total pembayaran pada
sistem yang dirancang ditentukan berdasarkan jenis dan jumlah lauk, bukan
berdasarkan luas area lauk pada citra.

Salah satu model yang digunakan dalam pengembangan \textit{object detection} adalah YOLO (\textit{You Only Look Once}). YOLO dirancang untuk mendeteksi objek dalam citra dengan melakukan lokalisasi dan klasifikasi dalam satu proses inferensi. Seiring perkembangannya, keluarga YOLO tidak hanya digunakan untuk \textit{object detection}, tetapi juga mendukung beberapa tugas \textit{Computer Vision} lainnya, termasuk \textit{image classification} dan \textit{instance segmentation} \cite{khanam2024}.

YOLO11 merupakan salah satu versi dari keluarga YOLO yang dikembangkan dengan peningkatan pada arsitektur dan efisiensi model. YOLO11 mendukung beberapa tugas \textit{Computer Vision}, termasuk \textit{instance segmentation}. Untuk tugas tersebut, digunakan varian YOLO11-seg yang menghasilkan informasi kelas dan \textit{pixel mask} untuk setiap objek yang terdeteksi \cite{khanam2024,qiu2025}. Output tersebut memungkinkan sistem untuk menentukan batas setiap lauk dan membedakan beberapa lauk sebagai instans yang terpisah sehingga jumlahnya dapat dihitung secara otomatis. Berdasarkan kebutuhan tersebut, YOLO11-seg digunakan sebagai model utama dalam penelitian ini.

Berdasarkan uraian tersebut, penelitian ini merancang \textit{proof-of-concept} sistem \textit{checkout} otomatis yang menyediakan opsi \textit{self-service} bagi pelanggan \textit{dine-in} restoran Nasi Padang. Pelanggan dapat mengambil lauk secara mandiri, kemudian sistem menggunakan YOLO11-seg untuk mendeteksi, melakukan \textit{instance segmentation}, dan menghitung setiap instans lauk pada piring. Hasil tersebut selanjutnya digunakan untuk menghitung total pembayaran berdasarkan harga tetap setiap kelas lauk. Penelitian ini dibatasi pada sembilan kelas, yaitu nasi, rendang, ayam goreng, perkedel, sayur nangka, daun singkong, telur dadar, telur kari, dan sambal.

Berdasarkan latar belakang tersebut, penelitian ini mengusulkan judul ``Deteksi Lauk Nasi Padang untuk Sistem \textit{Checkout} Otomatis dengan \textit{Instance Segmentation} Menggunakan YOLO11''.


\section{Rumusan Rancangan}

Berdasarkan latar belakang yang telah diuraikan, penelitian ini berfokus pada perancangan sistem \textit{checkout} otomatis yang memanfaatkan \textit{instance segmentation} untuk mengenali dan menghitung lauk Nasi Padang. Rumusan rancangan dalam penelitian ini adalah sebagai berikut:

\begin{enumerate}
    \item Bagaimana merancang dan melatih model \textit{instance segmentation} menggunakan arsitektur YOLO11-seg agar mampu mendeteksi dan melakukan segmentasi terhadap sembilan kelas lauk Nasi Padang serta menghasilkan prediksi setiap instans objek secara individual?

    \item Bagaimana mengevaluasi performa YOLO11-seg dalam mendeteksi, melakukan segmentasi, dan menghitung instans lauk Nasi Padang berdasarkan hasil pengujian model?

    \item Bagaimana merancang alur integrasi hasil deteksi dan \textit{instance counting} dari YOLO11-seg dengan sistem penetapan harga berbasis \textit{fixed price} untuk menghasilkan total tagihan secara otomatis?
\end{enumerate}


\section{Komponen Perancangan}

Sistem \textit{checkout} otomatis yang dirancang dalam penelitian ini terdiri atas tiga komponen utama yang saling terhubung dalam proses \textit{checkout}, yaitu kamera, model YOLO11-seg, dan program kalkulasi harga.

\begin{enumerate}
    \item \textbf{Kamera}

    Kamera berfungsi sebagai komponen akuisisi citra pada proses \textit{checkout}. Kamera digunakan untuk memperoleh citra piring yang telah berisi lauk pilihan pelanggan. Citra tersebut kemudian digunakan sebagai masukan bagi model YOLO11-seg untuk melakukan proses pengenalan lauk.

    \item \textbf{Model YOLO11-seg}

    YOLO11-seg merupakan komponen utama dalam pemrosesan \textit{Computer Vision} pada sistem. Model ini digunakan untuk mendeteksi dan melakukan \textit{instance segmentation} terhadap lauk yang terdapat pada citra piring. Dari hasil prediksi tersebut, sistem memperoleh kelas dan setiap instans lauk yang terdeteksi sehingga jumlah lauk dari masing-masing kelas dapat dihitung.

    \item \textbf{Program Kalkulasi Harga}

    Program kalkulasi harga berfungsi mengolah hasil pengenalan dan penghitungan instans dari YOLO11-seg menjadi total pembayaran. Program menggunakan harga tetap untuk setiap kelas lauk, kemudian menghitung total harga berdasarkan jenis dan jumlah instans lauk yang terdeteksi.
\end{enumerate}


\section{Spesifikasi Perancangan}

Spesifikasi perancangan sistem \textit{checkout} otomatis lauk Nasi Padang dalam penelitian ini terdiri atas tiga modul utama yang saling terhubung, yaitu modul akuisisi citra, modul pengenalan lauk menggunakan YOLO11-seg, dan modul kalkulasi harga.

\begin{enumerate}
    \item \textbf{Modul Akuisisi Citra}

    Modul akuisisi citra berfungsi memperoleh citra piring yang berisi lauk pilihan pelanggan menggunakan kamera. Citra yang diperoleh digunakan sebagai masukan bagi model YOLO11-seg. Sistem menggunakan citra digital dua dimensi dalam format RGB sebagai data masukan.

    \item \textbf{Modul Pengenalan Lauk dengan YOLO11-seg}

    Modul pengenalan lauk menggunakan YOLO11-seg untuk mendeteksi dan melakukan \textit{instance segmentation} terhadap lauk yang terdapat pada citra. Model menghasilkan prediksi kelas serta \textit{mask} untuk setiap instans objek yang terdeteksi.

    Hasil prediksi tersebut digunakan untuk menentukan jumlah instans pada setiap kelas. Setiap instans yang terdeteksi dihitung secara individual sehingga beberapa objek dari kelas yang sama tetap dapat dibedakan dan dihitung sebagai objek yang terpisah.

    Pengenalan dibatasi pada sembilan kelas, yaitu nasi, rendang, ayam goreng, perkedel, sayur nangka, daun singkong, telur dadar, telur kari, dan sambal.

    \item \textbf{Modul Kalkulasi Harga}

    Modul kalkulasi harga menggunakan jenis dan jumlah instans lauk yang telah dikenali untuk menentukan total pembayaran. Setiap kelas lauk memiliki harga tetap, kemudian total pembayaran dihitung berdasarkan jumlah instans dari masing-masing kelas. Penentuan harga tidak menggunakan berat, volume, luas \textit{pixel mask}, maupun ukuran porsi.
\end{enumerate}

\section{Kegunaan Rancangan}

Hasil perancangan sistem checkout otomatis ini diharapkan dapat memberikan kegunaan secara operasional maupun teknis dan akademis. Secara keseluruhan, rancangan ini menjadi \textit{proof-of-concept} penerapan \textit{self-service} dan otomatisasi proses \textit{checkout} pada restoran Nasi Padang dengan memanfaatkan \textit{Computer Vision}.

Secara operasional, sistem yang dirancang dapat digunakan sebagai alternatif mekanisme \textit{checkout} bagi pelanggan \textit{dine-in}. Sistem memanfaatkan YOLO11-seg untuk mengenali dan menghitung lauk yang terdapat pada piring, kemudian hasil tersebut digunakan untuk menentukan total pembayaran secara otomatis. Penerapan mekanisme ini diharapkan dapat mengurangi ketergantungan terhadap staf dalam proses pencatatan lauk pelanggan \textit{dine-in} serta memberikan alternatif pelayanan yang lebih mandiri.

Secara teknis dan akademis, penelitian ini memberikan penerapan dan evaluasi YOLO11-seg pada \textit{use-case} pengenalan dan penghitungan lauk Nasi Padang yang memiliki karakteristik objek dan kelas yang spesifik. Hasil evaluasi dapat memberikan gambaran mengenai kemampuan model dalam mendeteksi, melakukan \textit{instance segmentation}, dan menghitung setiap instans lauk. Dengan demikian, penelitian ini dapat menjadi referensi awal untuk penerapan \textit{Computer Vision} berbasis \textit{instance segmentation} pada sistem \textit{checkout} makanan dengan jenis dan jumlah objek sebagai dasar penentuan harga.

\section{Rancangan yang Sudah Dibuat}

Beberapa penelitian sebelumnya telah mengembangkan sistem dan metode yang berkaitan dengan pengenalan makanan, \textit{instance segmentation}, penghitungan instans, dan sistem \textit{checkout} otomatis. Beberapa rancangan yang relevan dengan penelitian ini adalah sebagai berikut:

\begin{enumerate}
    \item Perancangan yang berjudul ``Grab, Pay and Eat: Semantic Food Detection for Smart Restaurants'' oleh Aguilar et al. mengembangkan metode \textit{Semantic Food Detection} untuk menganalisis makanan pada nampan di lingkungan restoran. Metode tersebut menggabungkan proses lokalisasi, pengenalan, deteksi, dan segmentasi makanan untuk mendukung penerapan \textit{automatic billing} pada restoran \textit{self-service}. Penelitian tersebut memperoleh nilai F-measure sekitar 90\% pada dataset UNIMIB2016 \cite{aguilar2017}.

    \item Perancangan yang berjudul ``A Framework of Visual Checkout System Using Convolutional Neural Networks for Bento Buffet'' oleh Wu et al. mengembangkan sistem \textit{visual checkout} untuk restoran Bento Buffet menggunakan kamera Kinect dan CNN. Sistem melakukan pengenalan makanan serta menggunakan informasi kedalaman untuk memperkirakan volume makanan dan harga. Pendekatan pengenalan dua tahap yang digunakan menghasilkan akurasi sebesar 96,3\% pada dataset pengujian \cite{wu2021}.

    \item Perancangan yang berjudul ``A Real-Time Chinese Food Auto Billing System Based on Instance Segmentation'' oleh Guo et al. mengembangkan sistem perhitungan harga makanan Tiongkok secara otomatis menggunakan segmentasi. Sistem mengenali jenis piring berdasarkan bentuk dan ukuran, di mana setiap jenis piring merepresentasikan harga tertentu. Penelitian tersebut juga mengembangkan FoodSyn untuk menghasilkan data sintetis dan memperoleh nilai mIoU lebih dari 95\% pada pengujian yang dilakukan \cite{guo2023}.

    \item Perancangan yang berjudul ``SibNet: Food Instance Counting and Segmentation'' oleh Nguyen et al. mengembangkan jaringan multi-task untuk melakukan penghitungan dan segmentasi instans makanan secara bersamaan. SibNet dirancang untuk menangani karakteristik makanan dengan bentuk yang tidak beraturan serta kondisi objek yang saling menutupi melalui penggunaan \textit{instance seed} dan \textit{sibling relation network} \cite{nguyen2022sibnet}.

    \item Perancangan yang berjudul ``FoodMask: Real-time Food Instance Counting, Segmentation and Recognition'' oleh Nguyen et al. mengembangkan jaringan multi-task untuk melakukan penghitungan, segmentasi, dan pengenalan instans makanan secara bersamaan. FoodMask menghasilkan informasi kategori makanan dan instans pada tingkat piksel serta dirancang untuk melakukan pemrosesan secara \textit{real-time} pada citra yang dapat berisi beberapa jenis makanan \cite{nguyen2023foodmask}.

    \item Perancangan yang berjudul ``Indonesian Street Food Calorie Estimation Using Mask R-CNN and Multiple Linear Regression'' oleh Aditama dan Munir menerapkan Mask R-CNN untuk melakukan \textit{amodal instance segmentation} terhadap enam kelas jajanan Indonesia. Penelitian tersebut secara khusus mengevaluasi kondisi oklusi dan non-oklusi, dengan F1-score sebesar 0,821 pada kondisi oklusi dan 0,994 pada kondisi non-oklusi \cite{aditama2022}.
\end{enumerate}
